\chapter{光的量子性}

\problem{7.1}{在深度远大于表面波波长的液体中，表面波的传播速度满足如下规律：
\[v=\sqrt{\frac{\lambda}{2\pi}\left(g+\frac{4\pi^2 F}{\lambda^2\rho}\right)}\]
式中 $g$ 为重力加速度，$\rho$ 为液体密度，$F$ 为表面张力，$\lambda$ 为表面波的波长。试计算表面波的群速度。}
\answerarea

\problem{7.2}{测量二硫化碳的折射率的实验数据为：当 $\lambda'=589\,\mathrm{nm}$，$n'=1.629$；当 $\lambda''=656\,\mathrm{nm}$ 时，$n''=1.620$。试求波长 $589\,\mathrm{nm}$ 的光在二硫化碳中的相速、群速和群速折射率。}
\answerarea
\clearpage

\problem{7.3}{在测定光速的迈克耳孙旋转棱镜法中，设所用棱镜为正 $n$ 面棱柱体。试导出：根据棱镜的转速、反射镜距离等数据计算光速的公式。}
\answerarea

\problem{7.4}{试用光的相速度 $v$ 和 $\dfrac{\mathrm{d}v}{\mathrm{d}\lambda}$ 来表示群速度 $u=\dfrac{\mathrm{d}\omega}{\mathrm{d}k}$，再用 $v$ 和 $\dfrac{\mathrm{d}n}{\mathrm{d}\lambda}$ 表示群速度 $u=\dfrac{\mathrm{d}\omega}{\mathrm{d}k}$。}
\answerarea
\clearpage

\problem{7.5}{计算在下列各种色散介质中传播的各种不同性质的波的群速度：（1）$v=\text{常量}$（无色散介质，如空气中的声波）；（2）$v=a\sqrt{\lambda}$，$a$ 为常量（重力在水面上所引起的波）；（3）$v=\dfrac{a}{\sqrt{\lambda}}$（在水面上的表面张力波）；（4）$v=\dfrac{a}{\lambda}$（弹性薄片在弯曲时所产生的波）；（5）$v=\sqrt{c^2+b^2\lambda^2}$（电离层中的电磁波，其中 $c$ 是真空中的光速，$\lambda$ 是介质中的波长）；（6）$v=c\omega/\sqrt{\omega^2\varepsilon\mu-c^2a^2}$（在充满色散介质的直波导管中的电磁波，式中 $c$ 是真空中的光速，$a$ 是与波导管横截面的形状及大小有关的常量，$\varepsilon=\varepsilon(\omega)$ 是介质的电容率，$\mu=\mu(\omega)$ 是介质的磁导率）。}
\answerarea

\problem{7.6}{利用维恩公式求：辐射的最概然频率 $\nu_{\mathrm{m}}$、辐射的最概然波长 $\lambda_{\mathrm{m}}$、辐射的最大光谱密度 $(\varepsilon_\lambda)_{\mathrm{m}}$、辐射出射度 $M_0(T)$ 与温度 $T$ 的关系。}
\answerarea
\clearpage

\problem{7.7}{太阳光谱非常接近于 $\lambda_{\mathrm{m}}=480\,\mathrm{nm}$ 的绝对黑体的光谱。试求：（1）在 $1\,\mathrm{s}$ 内太阳由于辐射而损失的质量；（2）并估计太阳的质量减少 $1\%$（由于热辐射）所经历的时间。太阳的质量 $m_0$ 为 $2.0\times 10^{30}\,\mathrm{kg}$，太阳的半径 $r$ 是 $7.0\times 10^8\,\mathrm{m}$。}
\answerarea

\problem{7.8}{地球表面每平方厘米每分钟由于辐射而损失的能量平均值为 $0.546\,\mathrm{J}$。如有一个黑体，它在辐射相同的能量时，温度应为多少？}
\answerarea
\clearpage

\problem{7.9}{若有一黑体的辐出度等于 $5.70\,\mathrm{W/cm^2}$，试求与该辐射最大光谱强度相对应的波长。}
\answerarea

\problem{7.10}{用交流供电的灯丝温度是变动的。一电灯白炽时钨丝的平均温度为 $2\,300\,\mathrm{K}$，其最高和最低温度之差约为 $80\,\mathrm{K}$。问热辐射的总功率的最大值和最小值之比为多少？钨丝的辐射可当做黑体。}
\answerarea
\clearpage

\problem{7.11}{若将恒星表面的辐射近似地看做是黑体辐射，则可以用测量 $\lambda_{\max}$ 的办法来估计恒星表面的温度。现测得太阳的 $\lambda_{\max}$ 为 $510\,\mathrm{nm}$，北极星的 $\lambda_{\max}$ 为 $350\,\mathrm{nm}$，试求它们的表面温度。}
\answerarea

\problem{7.12}{小灯泡所耗的功率为 $1\,\mathrm{W}$，均匀地向各个方向辐射能量。设辐射的平均波长为 $500\,\mathrm{nm}$，试求在 $10\,\mathrm{km}$ 处每秒钟落在垂直于光线方向上每平方厘米面积上的光子数。}
\answerarea
\clearpage

\problem{7.13}{已知铯的逸出功为 $1.88\,\mathrm{eV}$，现用波长为 $300\,\mathrm{nm}$ 的紫外光照射，试求光电子的初动能和初速。}
\answerarea

\problem{7.14}{用波长为 $253\,\mathrm{nm}$ 的光照射钨的表面时，在光电管的电路中产生的光电流，由于外加 $1\,\mathrm{V}$ 的遏止电压而截止。已知钨的逸出功为 $4.5\,\mathrm{eV}$，试确定接触电势差。}
\answerarea
\clearpage

\problem{7.15}{波长为 $320\,\mathrm{nm}$ 的紫外线入射到逸出功为 $2.2\,\mathrm{eV}$ 的金属表面上，求光电子从金属表面逸出时的最大速度。若入射光的波长为原来的一半，出射光电子的最大动能是否增至两倍？}
\answerarea

\problem{7.16}{波长为 $0.1\,\mathrm{nm}$ 的 X 射线被碳块散射，在散射角为 $90^\circ$ 的方向上进行观察。试求（1）康普顿位移 $\Delta\lambda$；（2）反冲电子的动能。}
\answerarea
\clearpage

\problem{7.17}{已知入射光子的波长为 $0.003\,\mathrm{nm}$，而反冲电子的速度为光速的 $\beta$ 倍（$\beta=0.6$），试确定康普顿位移 $\Delta\lambda$。}
\answerarea

\problem{7.18}{在电子显微镜中，电子受到 $90\,\mathrm{kV}$ 的电压加速，如要观察到物质的分子结构（其大小为 $10^{-9}\,\mathrm{cm}$ 数量级），问显微镜的数值孔径应该多大？}
\answerarea
\clearpage

\sproblem{7.19}{（1）一只 $100\,\mathrm{W}$ 灯泡，$5\%$ 的功率辐射是可见光，假定可见光平均波长为 $500\,\mathrm{nm}$，则每秒可辐射的可见光光子的能量为多少？（2）假设灯泡为点光源，可以向空间各个方向发光，求在距离 $2\,\mathrm{m}$ 处每秒垂直通过单位面积的光子数。}
\answerarea

\sproblem{7.20}{人在黑暗中，眼睛的视网膜如能接收到波长为 $550\,\mathrm{nm}$ 的最大有效辐射为 $2\times 10^{-18}\,\mathrm{J}$，就能感知这一点光源。求眼睛这一观察阈值相当于多少个光子？}
\answerarea
\clearpage

\sproblem{7.21}{一电子束被电压为 $V$ 的电场加速。（1）求电子在被加速后的德布罗意波长；（2）求此德布罗意波的相速和群速；（3）把此德布罗意波射到一块单晶上，如入射方向与晶面成 $\theta$ 角，观察到散射波第一级强度极大值，求晶格常量 $d$。}
\answerarea

\sproblem{7.22}{在热核爆炸的火球中，测得瞬时温度为 $10^7\,\mathrm{K}$。（1）估算辐射最强的波长；（2）这种辐射最强波长的能量子是多少焦耳？}
\answerarea
\clearpage

\sproblem{7.23}{用关键词光镊（optical tweezer）或激光捕获（laser trapping），在互联网（如 Google 或 Yahoo）上查阅光镊技术在生物细胞检测和医学诊断方面的应用和进展，写一篇介绍报告。}
\answerarea
\clearpage
